\subsection{The cost of one user}

We first timed one user alone, back to back.
An AnySplat update takes 0.122~s of service time at $k=2$ frames, 0.395~s at $k=8$ and 0.859~s at $k=16$ (medians of 20 updates, one run per $k$). % G1
Its process holds 4,726--9,576~MiB, of which the weights take a constant 2,808~MiB. % G1
At $k=8$ and one update per user per second, this puts the compute bound at 2 users ($1/0.395~\text{s} = 2.5$ updates/s) and the memory bound at 5 ($49{,}140/8{,}640~\text{MiB} = 5.7$). % G1

A StreamSplat update, depth included, takes 0.190~s with \texttt{torch.compile} and 0.198~s without, of which the depth part takes 0.049~s, against the 0.09~s we had predicted. % G5
One user reaches 5.20 updates/s back to back, and its process holds 4,094~MiB. % G1, G5
The reconstruction model timed alone takes 0.152~s, matching the non-depth part, so the harness adds no measurable overhead. % G5

\subsection{Sharing mechanisms add no users}

In each of the four workloads of Table~\ref{tab:sharing}, time-slicing and MPS serve the same number of users in both repeats, and every failing point fails on latency. % G2, G3, G6, G7
Past the capacity the GPU completes about the same number of updates per second however it is shared, and every user's queue grows for the length of the run. % G2, G6
Up to the capacity, StreamSplat's worst p95 latency stays at 0.189--0.198~s for 1--4 users, close to one user's service time. % G6

MPS changes throughput by only a few percent in either direction, and its one measurable benefit in the AnySplat sweeps was in latency: with five users at $k=2$, where demand meets the GPU's rate, it cut the worst p95 from 0.318--0.330~s to 0.177--0.180~s. % G3, G6, G7
Time-slicing alone overlaps users more than we had predicted: four AnySplat users at $k=2$ run back to back reached 10.06 updates/s, above the 8.2 of one user's service time, and MPS reached 10.92. % G3
One possible reading, which we did not test, is that other processes' kernels fill the gaps a user leaves while it works on the CPU. % G3
For StreamSplat at one keyframe per second, the shared GPU completes 4.62--4.73 updates/s, 9--11\% below the 5.20 that one user reaches alone. % G5, G6

No sharing sweep ran out of memory: at the largest points run, four AnySplat users at $k=8$ peaked at 29,895--32,272~MiB and eight StreamSplat users at 32,831~MiB, of the card's 49,140~MiB. % G2, G6
At 4,094~MiB per StreamSplat user the card would hold about 12 users, three times the 4 that its compute serves. % G6

\begin{table}[t]
\centering
\footnotesize
\setlength{\tabcolsep}{3.5pt} % layout only
\begin{tabular}{@{}llcccc@{}}
\toprule
 & & \multicolumn{2}{c}{Users served} & \multicolumn{2}{c}{Saturated updates/s} \\ % layout only
\cmidrule(lr){3-4}\cmidrule(l){5-6} % layout only
Model & Rate & TS & MPS & TS & MPS \\
\midrule
AnySplat, $k=8$ & 1/s & 2 & 2 & 2.52--2.53 & 2.45--2.47 \\ % G2
AnySplat, $k=2$ & 2/s & 5 & 5 & 10.07--10.09 & 10.18--10.29 \\ % G3
StreamSplat & 1/s & 4 & 4 & 4.71--4.73 & 4.62--4.65 \\ % G6
StreamSplat & 3/s & 1 & 1 & 4.77--4.83 & 4.81--4.95 \\ % G7
\midrule
\multicolumn{2}{@{}l}{Batching server, batch size $b$} & 1 & 2 & 4 & 8 \\ % G9
\multicolumn{2}{@{}l}{Updates/s, back to back} & 4.92 & 4.83 & 4.70 & 4.52 \\ % G9
\bottomrule
\end{tabular}
\caption{Users served with one process per user under time-slicing (TS) and MPS, and the throughput once the GPU saturates, as the range over each sweep's failing points. % G2, G3, G6, G7
Rate is updates per user per second, each a keyframe for StreamSplat and $k$ frames for AnySplat.
Serving $N$ users requires both of 2 repeats to meet the SLO of Section~\ref{sec:setup}. % G2, G3, G6, G7
The StreamSplat sweep at 1/s skipped 5 users. % G6
The bottom rows give the batching server's back-to-back throughput, one run per batch size. % G9
}
\label{tab:sharing}
\end{table}

One server that runs several users' updates as one batch does not raise the rate either: back-to-back throughput falls as the batch grows (Table~\ref{tab:sharing}, bottom rows). % G9
A batch of 8 takes 8.7 times as long as one update (1.767 against 0.202~s), and the reconstruction model timed alone scales the same way, so the scaling belongs to the model on this GPU. % G9
GPU utilization rose from 72\% at batch 1 to 96\% at batch 8 while throughput fell, so the idle time at batch 1 was not capacity that batching could use. % G9
We had decided in advance to sweep the batching server only if batch 8 reached 1.2 times the batch-1 throughput; it reached 0.92 times and was never swept. % G9
We set aside one earlier batch-1 run of 4.32 updates/s, after seeing it, because another project's CPU-bound job had started 12~s before it; whether that job slowed the run was not tested. % G9
Whether one shared copy of the weights saves memory was not measured, since memory did not limit the per-user configurations. % G9

\subsection{Where an update's time goes}

Since neither sharing nor batching added users, we profiled one StreamSplat update to find work that could be removed. % G10
We split the forward pass, with bitwise identical output, into a single-frame stage (encoder, DINOv2 conditioning, static upsampler and head) and a pair stage (decoder, dynamic upsampler and head). % G10

At batch 1 in TF32, with each time below the median of 20 calls, the forward pass takes 158~ms and the depth stage, the depth model with its normalization, 55.1~ms. % G10
The released code encodes both keyframes of a pair at every update, so the single-frame stage of the previous keyframe, 34.0~ms or 21.5\% of the forward pass, repeats work done one update earlier. % G10
Under bf16 autocast the forward pass takes 66~ms, 2.4 times faster, and attention explains most of the difference. % G10
In FP32, PyTorch's scaled-dot-product attention has no FlashAttention kernel and runs a memory-efficient kernel at 13--19~TFLOP/s, while in bf16 FlashAttention runs at 142--173~TFLOP/s, and attention takes 67\% of kernel time in TF32 against 25\% in bf16. % G10
The GPU is busy for 86\% of a batch-1 forward pass in TF32 and 78\% in bf16, and one forward pass makes 14 stream synchronizations and launches 1,315 kernels in TF32. % G10
Batching raises the cost per sample mainly in the two upsamplers, whose time per sample at batch 4 is 1.46 and 1.25 times that at batch 1, for reasons we did not test. % G10

By thresholds fixed before the profile, we dropped two scheduling ideas: batching only the stages that benefit from it would save 2.9~ms per update at batch 4, against 10~ms required, and running the depth and pair stages on two CUDA streams was 0.1--0.6\% slower, against a required 5\% gain. % G10
By the same thresholds we kept the feature cache (above a 12\% share of model time), bf16 autocast (0.42 of the TF32 time against 0.87) and host-synchronization removal with CUDA graphs (GPU busy 86\% of the time against 90\%), which the rest of this section measures in the server. % G10; G11-G16
